\emph{Setting.} $\Cl(2,4)$ on orthonormal generators $e_{0}, \dots, e_{5}$ with squares $(+,-,-,-,+,-)$, $\nO = \tfrac12(e_{4} + e_{5})$, $\nInf = e_{4} - e_{5}$, so $\nO \cdot \nInf = 1$, both null, and $i_{6} = e_{0}e_{1}e_{2}e_{3}e_{4}e_{5}$ (the chart and orientation of Theorem \claimref{synthesis.frozen-point-integrand-on-tractor-transport-is-weyl-squared}). For a bivector $B$ and a vector $y$, $B \times y = \tfrac12(By - yB) = B \cdot y$. All statements are pointwise and algebraic in the fields at a point, so we work in an orthonormal frame of the derived metric with $\hb = 1$ there; the frame-covariant form follows by writing every slot in frame components. Four-forms: a two-form is $F = \tfrac12 F_{ef}e^{e} \wedge e^{f}$; the coefficient of $e^{0} \wedge e^{1} \wedge e^{2} \wedge e^{3}$ in a wedge of one-forms $\alpha_{1}, \dots, \alpha_{4}$ is $\sum_{p}\sgn(p)\alpha_{1,p_{0}}\cdots\alpha_{4,p_{3}}$, with the factors $\tfrac12$ per two-form; this fixes $L_{1}$, $L_{3}$, $L_{5}$ and $St$ as in the check.

\emph{(i) Broken components.} With $I = \nInf + \rho\nO$: $(\nInf \wedge a)\cdot\nInf = 0$, $(\nInf \wedge a)\cdot\nO = \nInf(a \cdot \nO) - a(\nInf \cdot \nO) = -a$, $(\nO \wedge s)\cdot\nInf = -s$, $(\nO \wedge s)\cdot\nO = 0$, and $(\nInf \wedge \nO)\cdot\nInf = -\nO$, $(\nInf \wedge \nO)\cdot\nO = \nInf$; a Lorentz bivector annihilates $\nO$ and $\nInf$. Hence $DI(a) = \rho a - S(a) + \beta(a)(\nInf - \rho\nO)$ and $DX(a) = a - \beta(a)\nO$ (verified on a generic slot and dilation at two values of $\rho$). That these vanish together exactly for $S = \rho\,\id$, $\beta = 0$ is the slot lemma of Theorem \claimref{synthesis.parallel-frozen-point-on-matter-regions-is-the-field-space-premise}(i); only the explicit form is used below.

\emph{The Lorentz block.} $d\Omega_{T}$ and the commutators of $\Omega$ and of $\beta\,\nInf \wedge \nO$ with the other parts have no component on spacetime bivectors, and $(\nInf \wedge a) \times (\nInf \wedge b) = 0 = (\nO \wedge s) \times (\nO \wedge t)$ since both vectors are null and orthogonal to spacetime. So the spacetime block of the field strength is $-\mathcal{R}(A \wedge B) + T_{S}(A, B)$ with $T_{S}$ the spacetime part of $\Theta(A) \times \Theta(B)$, $\Theta(A) = -\nInf \wedge A + \nO \wedge S(A)$, linear in $S$; for $S = \Sch$ it is the Kulkarni-Nomizu block of $\Sch$, so the block is $-\mathcal{W}$ (Theorem \claimref{synthesis.gtg-tractor-gauge-field-reproduces-tractor-connection}(iv)), and for $S = \rho\,\id$ it is $2\rho A \wedge B$ ((v) of the same theorem). The slot enters this block without derivatives (its derivative sits in the $\nO$ block), so every density below is algebraic in $S$. Since $I \wedge V = I \wedge X = \nInf \wedge \nO$ (Theorem \claimref{dark-energy.frozen-point-makes-stelle-west-vector}), every density here depends on its two-form and one-form arguments only through their spacetime parts: a component along $\nO$ or $\nInf$ wedges to zero against $\nInf \wedge \nO$. In particular $\beta$ drops out of $L_{1}$, $L_{3}$, $L_{5}$ and $St$.

\emph{(iii) Stiffness and the slot Hessian.} By (i) the spacetime parts are $DI = -(S - \rho\,\id)$ and $DX = \id$, so $St = \sum_{p}\sgn(p)\,\epsilon(I, X, M e_{p_{0}}, M e_{p_{1}}, e_{p_{2}}, e_{p_{3}})$ with $M = S - \rho\,\id$, a quadratic form in $M$ alone. The check evaluates it on thirty-nine slots, with and without a dilation, and finds $Q_{2}(M)$; since both sides are quadratic forms on a sixteen-dimensional space and the test set contains the identity shifts of all sixteen matrix units and of twenty pairs of them, the polarised identity is fixed. For $L_{1}(S) = \epsilon_{4}(-\mathcal{R} + T_{S}, -\mathcal{R} + T_{S})$ the quadratic part is $\epsilon_{4}(T_{S}, T_{S})$, which is the same $Q_{2}(S)$ (same test set), and the linear part $2\epsilon_{4}(-\mathcal{R}, T_{S})$ is bilinear in $(\mathcal{R}, S)$. On a basis of the thirty-six curvatures antisymmetric in each pair (no pair symmetry or Bianchi identity assumed, so torsion is allowed) and the sixteen matrix units the check verifies $2\epsilon_{4}(-\mathcal{R}, T_{S}) = -2Q_{2}(S, \Sch_{*})$, $Q_{2}(A, B) = 2(\tr A\,\tr B - \tr AB)$, with $\Sch_{*}$ the Schouten map of the contraction $R_{bd} = \sum_{a}\eta^{aa}R_{adab}$. Completing the square gives $L_{1}(S) = L_{1}(\Sch_{*}) + Q_{2}(S - \Sch_{*})$ for every curvature. $Q_{2}$ is nondegenerate: its Hessian is $M \mapsto 4((\tr M)\id - M^{T})$, invertible in dimension four. For torsion-free $\Omega$, $\Sch_{*}$ is the stored Schouten map and $L_{1}(\Sch_{*}) = \tfrac14 W_{abcd}W^{abcd}$ (Theorem \claimref{synthesis.frozen-point-integrand-on-tractor-transport-is-weyl-squared}(iii)); this is the conformal-gauge construction of conformal gravity read in the series (Fact \claimref{synthesis.imported-conformal-gravity-bach-equations-and-conformal-gauge-form}(b)); the elimination is Theorem \claimref{synthesis.frozen-point-integrand-gives-einstein-only-with-the-symmetric-slot-frozen}(ii) and is not claimed here, only the identity of its Hessian form with $St$.

\emph{(iv) Interpolation.} By (iii), $L_{1} + \mu\,St = L_{1}(\Sch_{*}) + Q_{2}(S - \Sch_{*}) + \mu\,Q_{2}(S - \rho\,\id)$. For $\mu \neq -1$ its Hessian is $(1 + \mu)$ times that of $Q_{2}$, nondegenerate, and the unique stationary slot is $S_{\mu} = (\Sch_{*} + \mu\rho\,\id)/(1 + \mu)$. There $S_{\mu} - \Sch_{*} = \tfrac{\mu}{1+\mu}(\rho\,\id - \Sch_{*})$ and $S_{\mu} - \rho\,\id = \tfrac{1}{1+\mu}(\Sch_{*} - \rho\,\id)$, so the value is $L_{1}(\Sch_{*}) + \bigl(\tfrac{\mu^{2}}{(1+\mu)^{2}} + \tfrac{\mu}{(1+\mu)^{2}}\bigr)Q_{2}(\rho\,\id - \Sch_{*}) = L_{1}(\Sch_{*}) + t\,Q_{2}(\rho\,\id - \Sch_{*})$. At $\mu \to \infty$ ($t = 1$) this is $L_{1}(\rho\,\id)$, the frozen-point density, which the check evaluates as $\tfrac14(R_{abcd}R^{abcd} - 4R_{ab}R^{ab} + R^{2}) - \tfrac{\Lambda}{3}(R - 2\Lambda)$ on a torsion-free sample at two values of $\rho$ (the standard expansion of the MacDowell-Mansouri density into Euler, Einstein-Hilbert and volume terms, Fact \claimref{dark-energy.imported-macdowell-mansouri-theorem-stelle-west} as used in Theorem \claimref{synthesis.macdowell-mansouri-principle-fixes-the-trace-free-part}); hence the displayed density $(1 - t)L_{1}(\Sch_{*}) + t\,L_{1}(\rho\,\id)$. In it the Einstein-Hilbert coefficient is $-t\Lambda/3$ and the Weyl-square coefficient $(1-t)/4$ (the Gauss-Bonnet part is topological), a ratio of magnitude $3(1-t)/(4t\Lambda) = 3/(4\mu\Lambda)$. $Q_{2}$ is indefinite ($Q_{2}(\mathrm{diag}(m)) = 4\sum_{i<j}m_{i}m_{j}$), so $St$ is a penalty on the auxiliary slot, not a mass term. The check solves the sixteen linear slot equations exactly for $\mu \in \{0, 1, 3, -\tfrac12\}$ at both values of $\rho$ (Hessian and gradient computed from the densities themselves, not from the formula) and compares the solution and the value.

\emph{(v) Frame and slot.} $T_{S}(A, B)$ is the spacetime part of $(\nInf \wedge A)\times(\nO \wedge S(B)) + (\nO \wedge S(A))\times(\nInf \wedge B)$ up to sign, which is symmetric under exchanging the frame and the slot. With $\Omega$ fixed (its curvature has fixed coordinate components) and $S = \rho H$ for the frame $H = \hb^{-1}$, a frame variation $\delta H$ changes the block by $T(\delta H, \rho H) + T(H, \rho\,\delta H) = 2\rho\,T(H, \delta H)$, which is $2\rho$ times the change under the slot variation $\delta S = \delta H$. So the frame derivative is $2\rho$ times the slot derivative, the latter being $2(\mathcal{G} + \Lambda\,\id)$ at the frozen value (both computed exactly by a five-point stencil, exact for polynomials of degree at most four, at two values of $\rho$). Dividing by $Q(I) = 2\rho$, the prefactor of the frozen-point functional, gives the stored frame equation $2(\mathcal{G}(a) + \Lambda a)$.

\emph{(ii) The Stelle-West family.} At $S = \rho\,\id$ the slot derivatives are affine in the curvature and Lorentz-equivariant. For torsion-free $\Omega$ the equivariant linear maps from algebraic curvature tensors to linear maps of the spacetime vectors are spanned by the trace-free Ricci map and the scalar curvature times the identity (the Weyl part has no such image), so each derivative is $a\,\Ric_{0} + b\,R\,\id + c\,\id$ and is fixed by one sample with nonzero trace-free Ricci and scalar curvature together with the flat sample. The check computes the derivatives of $L_{1}$, $L_{3}$, $L_{5}$ exactly at $\rho \in \{\tfrac37, -\tfrac25\}$ on such a sample and on flat curvature and finds $2(\mathcal{G} + \Lambda\,\id)$, $(\mathcal{G} + 2\Lambda\,\id)/\rho$ and $(12/\rho)\id$ ($L_{5}$ has no curvature). The combination gives the displayed slot equation. Its vanishing for every torsion-free curvature requires $2c_{1} + c_{3}/\rho = 0$ and $2c_{1}\Lambda + 2c_{3}\Lambda/\rho + 12c_{5}/\rho = 0$, that is $c_{3} = -2\rho c_{1}$ and $c_{5} = \rho\Lambda c_{1}/6 = \rho^{2}c_{1}$. On that line the check finds the slot derivative zero at the frozen value also for a curvature without pair symmetry, and the frozen value of the density equal to $c_{1}\epsilon(I, V, \mathcal{R}, \mathcal{R})$, the Euler form of $\Omega$ (on the torsion-free sample $\tfrac14$ of the Gauss-Bonnet density), a characteristic form whose variation vanishes identically; at a generic slot the density differs, and a perturbed coupling ($c_{5} = 2\rho^{2}c_{1}$) has a nonzero slot derivative. The family is that of Theorem \claimref{synthesis.stelle-west-norm-is-no-local-output} with $V = X - I/(2\rho)$; there $\Omega_{T}$ was reduced, so $DV$ was the frame; here $DV = DX - DI/(2\rho)$ also carries the broken components.

\emph{Curvature-independent slot potentials.} If $U(S)$ does not involve the curvature, the slot derivative of $c_{1}L_{1} + U$ at $\rho\,\id$ is $2c_{1}(\mathcal{G} + \Lambda\,\id) + U'(\rho\,\id)$ with $U'(\rho\,\id)$ a fixed map; it vanishes for every curvature only if $c_{1} = 0$. A multiplier term $\langle \lambda \wedge DI\rangle$ with a vector-valued three-form $\lambda$ (polynomial, no inverse frame) has $DI = 0$ as its $\lambda$ equation at any coupling; it writes the premise into the action.

\emph{Gauge.} With $I$ and $X$ fixed, the densities are invariant under their joint stabiliser, the Lorentz rotations; a special-conformal rotor of the parabolic group moves $I$, so the constant-$I$ description is the unitary gauge of the Higgs reading.

\emph{Reading.} Statements (ii) to (iv) concern only which field space is varied; the integrand is the stored one throughout. The Einstein condition is never assumed: it appears as the condition under which the frozen slot is stationary ((ii), (v)). Stelle and West (1980, cited from memory) break the de Sitter group to Lorentz with a vector whose covariant derivative is the frame; the analogue for the first step would need the covariant derivative of $I$ to vanish, which the computation shows no finite stiffness enforces; the field-space premise reads, equivalently, as infinite stiffness. The fourth-order character of an added Weyl square (Stelle 1977, cited from memory) is mentioned in the statement's reading only.
